\documentclass[12pt,a4paper,UTF8]{ctexart}
\usepackage[a4paper,top=22mm,bottom=22mm,left=22mm,right=22mm,headheight=15pt]{geometry}
\usepackage{amsmath,amssymb,mathtools,bm}
\usepackage{graphicx,xcolor,booktabs,longtable,array,enumitem,fancyhdr,caption,float,microtype,setspace}
\newcommand{\Needspace}[1]{\par\begingroup\dimen0=#1\relax\dimen2=\pagegoal\advance\dimen2 by-\pagetotal\ifdim\dimen0>\dimen2\newpage\fi\endgroup}
\newcommand{\needspace}[1]{\par\begingroup\skip0=\lastskip\vskip-\skip0\dimen0=#1\relax\advance\dimen0 by\skip0\vskip\dimen0\penalty-100\vskip-\dimen0\vskip\skip0\penalty9999\endgroup}
\usepackage[hidelinks,unicode]{hyperref}
\usepackage[most]{tcolorbox}
\definecolor{ContestBlue}{HTML}{17365D}
\hypersetup{pdftitle={第 045 套},pdfauthor={}}
\setstretch{1.24}
\setlength{\parindent}{2em}
\setlength{\parskip}{0.35em}
\setlength{\emergencystretch}{2em}
\setlist{itemsep=0.35em,topsep=0.4em}
\captionsetup{font=small,labelfont=bf,justification=centering}
\setcounter{secnumdepth}{0}
\pagestyle{fancy}
\fancyhf{}
\fancyhead[L]{\small 第 045 套}
\fancyfoot[C]{\thepage}
\renewcommand{\headrulewidth}{0.4pt}
\ctexset{
  section={format=\Large\bfseries\color{ContestBlue},beforeskip=1.2em,afterskip=0.7em},
  subsection={format=\large\bfseries\color{ContestBlue},beforeskip=1.0em,afterskip=0.5em},
  subsubsection={format=\normalsize\bfseries,beforeskip=0.8em,afterskip=0.35em}
}
\makeatletter
\newcommand{\tightlist}{\setlength{\itemsep}{0pt}\setlength{\parskip}{0pt}}
\newcommand{\maxwidth}{\ifdim\Gin@nat@width>\linewidth\linewidth\else\Gin@nat@width\fi}
\newcommand{\maxheight}{\ifdim\Gin@nat@height>0.72\textheight 0.72\textheight\else\Gin@nat@height\fi}
\newcommand{\pandocbounded}[1]{\begingroup\setkeys{Gin}{width=\maxwidth,height=\maxheight,keepaspectratio}#1\endgroup}
\makeatother
\begin{document}
\begin{titlepage}
  \centering
  \vspace*{58mm}
  {\Huge\bfseries\color{ContestBlue} 第 045 套\par}
  \vfill
  \vspace*{28mm}
\end{titlepage}
\section{第 1 题}\label{ux7b2c-1-ux9898}

三体星系中的居民可利用径向速度法探知行星：恒星在行星引力作用下绕共同质心运动，使其光谱产生周期性多普勒频移。估算太阳受地球影响而产生的往返运动最大速率
\(v\)。已知地球质量 \(M_E=6.0\times10^{24}\,\mathrm{kg}\)、太阳质量
\(M_S=2.0\times10^{30}\,\mathrm{kg}\)、万有引力常量
\(G=6.67\times10^{-11}\,\mathrm{N\,m^2\,kg^{-2}}\)，地球公转周期取一年。

\begin{figure}[H]
\centering
\includegraphics[width=0.564\linewidth,height=\textheight,keepaspectratio,alt={题1　太阳与地球绕共同质心运动}]{../figures/question-01.pdf}
\caption{题1　太阳与地球绕共同质心运动}
\end{figure}

\newpage

\section{第 2 题}\label{ux7b2c-2-ux9898}

质量为 \(M\)、半径为 \(R\) 的均匀圆柱体置于倾角为 \(\theta\)
的斜面上，可沿斜面纯滚动而不打滑。斜面在地面参考系中以水平加速度 \(a_0\)
从斜面高端指向低端。求圆柱体在地面参考系中的加速度。

\begin{figure}[H]
\centering
\includegraphics[width=0.746\linewidth,height=\textheight,keepaspectratio,alt={题2　水平加速斜面上的纯滚动圆柱}]{../figures/question-02.pdf}
\caption{题2　水平加速斜面上的纯滚动圆柱}
\end{figure}

\newpage

\section{第 3 题}\label{ux7b2c-3-ux9898}

一个固定磁矩为 \(m_0\) 的小磁铁竖直放置，其正上方有磁矩为 \(m\)、质量为
\(M\) 的小磁铁，两磁矩反向。固定磁铁远场轴线上距其 \(h\)
处的磁感应强度为 \[
B(h)=\frac{\mu_0}{4\pi}\frac{2m_0}{h^3}.
\]

（1）求上方小磁铁的平衡高度 \(h\)。

（2）求其在竖直方向作小振动的频率 \(f\)。

（3）为稳定角度方向，令上方磁铁为均匀球体、过质心转动惯量为
\(I_0\)，磁矩相对竖直方向偏转小角 \(\alpha\)，并以角速度 \(\omega_0\)
绕磁矩方向自转。求稳定进动时绕竖直轴的角速度 \(\Omega\)。

（4）限制小磁铁质心不动，只考虑角度运动。在原竖直场 \(B_0\)
上叠加竖直方向高频弱交变场
\(B_1\cos\omega t\)，求高频近似下稳定所需的最小振幅 \(B_{1\min}\)。

\begin{figure}[H]
\centering
\includegraphics[width=0.262\linewidth,height=\textheight,keepaspectratio,alt={题3　反向磁矩的悬浮、进动与高频稳定}]{../figures/question-03.pdf}
\caption{题3　反向磁矩的悬浮、进动与高频稳定}
\end{figure}

\newpage

\section{第 4 题}\label{ux7b2c-4-ux9898}

一根质量可忽略、长为 \(2c\) 的刚杆两端为 \(A,B\)，分别固定质量
\(m_1,m_2\)，且 \(m_1<m_2\)。一根长为 \(2a\)、不可伸长的轻绳连接
\(A,B\)，并跨过固定光滑销钉 \(C\)，其中 \(a>c\)。三角形 \(ABC\)
始终在竖直平面内。

（1）求第 I 种竖直平衡时系统质心到 \(C\) 的距离 \(l_I\)。

（2）求第 II 种竖直平衡时系统质心到 \(C\) 的距离 \(l_{II}\)。

（3）若 \(\lambda m_2<m_1<m_2\) 时还存在第 III 种平衡，求 \(\lambda\)。

（4）求第 III 种平衡时质心到 \(C\) 的距离 \(l_{III}\)。

（5）在第 III 种平衡存在时，判断第 I 种平衡的稳定性。

（6）在第 III 种平衡存在时，判断第 II 种平衡的稳定性。

（7）当 \(m_1<\lambda m_2\) 时，判断第 I 种平衡的稳定性。

（8）取
\(a=0.696\,\mathrm m\)、\(c=0.313\,\mathrm m\)、\(m_1=3.20\,\mathrm{kg}\)、\(m_2=3.55\,\mathrm{kg}\)、\(g=9.8\,\mathrm{m\,s^{-2}}\)，求第
II 种平衡附近的小振动角频率。

\begin{figure}[H]
\centering
\includegraphics[width=0.8\linewidth,height=\textheight,keepaspectratio,alt={题4　刚杆、两端质点与跨钉轻绳的三类平衡}]{../figures/question-04.pdf}
\caption{题4　刚杆、两端质点与跨钉轻绳的三类平衡}
\end{figure}

\newpage

\section{第 5 题}\label{ux7b2c-5-ux9898}

一根无限长导体圆柱与无限大接地导体平面平行。圆柱轴线到平面的距离为
\(l\)，圆柱半径为 \(r\)。

（1）在 \(r/l\ll1\) 时，求系统单位长度电容，并给出小参数展开。

（2）不再假设 \(r/l\) 很小，求系统单位长度电容的精确表达式。

\begin{figure}[H]
\centering
\includegraphics[width=0.644\linewidth,height=\textheight,keepaspectratio,alt={题5　接地平面旁的无限长导体圆柱}]{../figures/question-05.pdf}
\caption{题5　接地平面旁的无限长导体圆柱}
\end{figure}

\newpage

\section{第 6 题}\label{ux7b2c-6-ux9898}

分子 \(A\) 可发生二聚反应 \[
2A\rightleftharpoons A_2.
\] 一分子 \(A\) 与一分子 \(A_2\) 在温度 \(T\) 时的平均分子总动能均为
\(3kT\)，均视为理想气体，并忽略 \(A_2\) 与两个 \(A\)
的内势能差。温度低于 \(T_0\) 时全部为 \(A_2\)；在 \(T_0<T<2T_0\)
时，\(A_2\) 的解离比例为 \[
\lambda=\frac{T-T_0}{T_0};
\] 温度高于 \(2T_0\) 时全部解离。

（1）在 \(T_0<T<2T_0\) 时求状态方程，以完全未解离时 \(A_2\) 的物质的量
\(n_0\) 表示。

（2）求这一温区混合气体的定容摩尔热容 \(C_V\) 与定压摩尔热容 \(C_P\)。

（3）求这一温区的绝热方程。

（4）初态为 \(V_0,T_0\)。依次经历：在 \(T_0\) 等温膨胀到
\(2V_0\)；绝热压缩升温到 \(2T_0\)；在 \(2T_0\)
等温压缩；绝热膨胀回到初态。求低温等温过程吸热 \(Q_1\)、高温等温过程放热
\(Q_2\)，并讨论循环性能。

\begin{figure}[H]
\centering
\includegraphics[width=0.8\linewidth,height=\textheight,keepaspectratio,alt={题6　二聚---解离气体的四过程循环}]{../figures/question-06.pdf}
\caption{题6　二聚---解离气体的四过程循环}
\end{figure}

\newpage

\section{第 7 题}\label{ux7b2c-7-ux9898}

肥皂泡膜有两个表面，每个表面的表面张力系数为
\(\sigma\)，膜的面质量密度为 \(\rho\)。

（1）推导二维无限肥皂膜的小振幅横波方程，并求波速。

（2）半径为 \(r\) 的水平圆形铁丝框覆盖肥皂膜，重力加速度 \(g\)
沿对称轴方向，且 \(\rho\) 足够小。求距圆心径向距离 \(x\) 处膜的下降高度
\(h(x)\)，特别求中心下降 \(h_0\)。

（3）两同轴圆环半径分别为 \(r_1>r_2\)，轴向间距为
\(h_0\)，其间张有肥皂膜。取大圆环平面为
\(z=0\)，忽略斜率平方以上的小量，求重力场中轴对称膜面 \(z(x)\) 的方程。

\begin{figure}[H]
\centering
\includegraphics[width=0.86\linewidth,height=\textheight,keepaspectratio,alt={题7　题设与符号}]{../figures/question-07.pdf}
\caption{题7　题设与符号}
\end{figure}

\newpage

\section{第 8 题}\label{ux7b2c-8-ux9898}

相对论换系时波的相位不变： \[
\phi(x,y,z,t)=\phi'(x',y',z',t').
\] 介质折射率 \(n\) 定义在介质静止参考系中。

（1）一团折射率为 \(n\) 的介质以速度 \(\beta c\) 沿 \(x\)
轴负方向相对地面运动。地面系中一束真空波长为 \(\lambda\) 的光沿 \(+x\)
方向入射该介质，地面系测得介质内波长为 \(\xi\lambda\)，求 \(\xi\)。

（2）一束平面波沿 \(+x\) 方向传播。折射率为 \(n\) 的介质沿 \(+x\)
方向流动，速度随离 \(x\) 轴距离 \(\rho\) 变化： \[
u(\rho)=\frac{\rho^2}{\rho_0^2}c,
\qquad \rho\ll\rho_0.
\] 平面波传播距离 \(f\) 后聚焦，求 \(f/\rho_0\)。

\begin{figure}[H]
\centering
\includegraphics[width=0.8\linewidth,height=\textheight,keepaspectratio,alt={题8　运动介质中的波矢变换与渐变流聚焦}]{../figures/question-08.pdf}
\caption{题8　运动介质中的波矢变换与渐变流聚焦}
\end{figure}
\end{document}
